% !TeX program = xelatex
% Edit this file directly; cv-style.sty controls the layout.
\documentclass[10pt,a4paper]{article}
\usepackage{cv-style}
% Shared contact details and content review date for both languages.
% Change the date when the content is reviewed, not just when recompiling.
\newcommand{\CVUpdated}{2026-09-06}
\newcommand{\CVEmail}{23300240019@m.fudan.edu.cn}
\newcommand{\CVPhone}{+86 158-0701-0023}
\newcommand{\CVWebsite}{https://16yunh.github.io/}
\newcommand{\CVWebsiteLabel}{16yunh.github.io}
\newcommand{\CVGitHub}{https://github.com/16yunH}
\newcommand{\CVGitHubLabel}{github.com/16yunH}
\newcommand{\CVPortrait}{assets/portrait.jpg}

\hypersetup{pdftitle={洪运 - 中文简历},pdflang={zh-CN}}
\begin{document}
\cvsetup{zh}
\cvheader{洪运}{计算机科学与技术 \textbar{} 复旦大学}
在复旦大学 MagicLab、丁文超老师指导下持续参与科研，目前研究长时序机器人操作中世界动作模型的执行反馈与子任务交接。

\cvsection{教育经历}
\begin{cventry}{复旦大学}{2023.09 - 至今}{计算机科学与技术 本科在读 \textbar{} 上海}
\end{cventry}

\begin{cventry}{得克萨斯大学奥斯汀分校}{2025.08 - 2025.12}{计算机科学 本科交换 \textbar{} 美国奥斯汀}
\begin{cvitems}
\item GPA：4.0/4.0，共 13 学分；机器学习、计算机视觉、操作系统、概率与随机过程四门课程均为 A。
\end{cvitems}
\end{cventry}

\cvsection{科研经历}
\begin{cventry}{面向机器人操作的世界动作模型研究}{当前研究}{复旦大学 MagicLab \textbar{} 指导老师：丁文超}
\begin{cvitems}
\item 基于 LingBot-VA 与 LIBERO-Long 搭建执行反馈框架，研究如何将任务进度与预测未来偏差信号作为条件反馈给同一策略，修正后续动作。
\item 通过状态重放与受控干预，区分模型输出变化和实际任务恢复，分析子任务交接时机械臂与夹爪构型、可达性及流式观测历史的作用。
\item 探索如何通过安全、可执行的动作到达有利的仿真起点，再由世界动作模型继续完成操作任务。
\end{cvitems}
\end{cventry}

\begin{cventry}{部分可观测环境下的自动驾驶风险地图学习}{2024.10 - 2025.05}{复旦大学 MagicLab \textbar{} 指导老师：丁文超}
\begin{cvitems}
\item 协助师兄完成 Transformer 风险预测框架搭建，以及模型训练、调试与结果验证。
\item 相关论文发表于 IEEE RA-L 2026，本人为第四作者，完整书目信息见下。
\end{cvitems}
\end{cventry}

\begin{cventry}{多任务机器人的主动感知与信息获取}{2025 年秋季}{UT Austin RobIn Lab \textbar{} 指导老师：Junhong Xu}
\begin{cvitems}
\item 搭建包含未知物体属性与未知空间位置的多任务、部分可观测机器人仿真 benchmark。
\item 搭建 VLA 基线训练与评测流程，用于研究模仿学习、强化学习、模仿预训练后强化微调三条路线。
\item 调研机器人操作探索与长时序 POMDP 文献；项目主要完成 benchmark 与基线评测流程。
\end{cvitems}
\end{cventry}

\begin{cventry}{大模型推理反思与自查机制}{2024.09 - 2025.06}{复旦大学自然语言处理组 \textbar{} 指导老师：黄萱菁}
\begin{cvitems}
\item 基于 Llama3.1-8B-Instruct 构建 Self-Critic 数据，采用树搜索与逆向推理生成多样化纠错轨迹。
\item 搭建数据合成与训练流程，在 GSM8K 上评估推理与反思方法。
\end{cvitems}
\end{cventry}

\cvsection{论文发表}
% Shared publication record; maintain author order here.
{\cvsmall
\textbf{Learning a Unified Risk Map for Autonomous Driving in Partially Observable Environments}\par
Jie Jia, Yaofeng Su, Zeyu Bao, \textbf{Yun Hong}, Bingzhao Gao, Zhongxue Gan, Wenchao Ding\par
\textit{IEEE Robotics and Automation Letters}, 11(4): 4457-4464, 2026.\par
\href{https://doi.org/10.1109/LRA.2026.3666393}{doi:10.1109/LRA.2026.3666393}\par}


\newpage
\cvsection{实习经历}
\begin{cventry}{字节跳动}{2026.01 - 2026.06}{测试开发实习生 \textbar{} 抖音研发 · 直播流媒体质量保障 \textbar{} 上海}
\begin{cvitems}
\item 开发三类 AI 提效工具，分别面向日常监测、测试用例风险检查和报警汇总分析。
\item 重构并新增巡检测试用例；在既定评测范围内，监测准确率由 89\% 提升至 100\%，覆盖率由 95\% 提升至 98\%。
\end{cvitems}
\end{cventry}

\cvsection{项目经历}
\begin{cventry}{LCDP-Sim：语言条件扩散策略机械臂控制}{始于 2025.12}{\href{https://github.com/16yunH/LCDP-Sim}{https://github.com/16yunH/LCDP-Sim}}
\begin{cvitems}
\item 构建结合 CLIP 与 U-Net DDPM/DDIM 扩散策略的视觉语言动作系统，实现 16 步 Action Chunking。
\item 在 ManiSkill2 中搭建数据采集、训练与闭环评测流程。
\end{cvitems}
\end{cventry}

\begin{cventry}{MapNavigation：路径规划与导航系统}{始于 2024.12}{\href{https://github.com/16yunH/MapNavigation}{https://github.com/16yunH/MapNavigation}}
\begin{cvitems}
\item 集成 OpenStreetMap 与高德地图 API，实现路径规划、地点搜索、混合选点及多道路类型速度约束。
\end{cvitems}
\end{cventry}

\begin{cventry}{NeuralStyle：模块化神经风格迁移工具箱}{始于 2025.06}{\href{https://github.com/16yunH/NeuralStyle}{https://github.com/16yunH/NeuralStyle}}
\begin{cvitems}
\item 开发模块化 Python 风格迁移工具，支持多图批处理、可配置实验流程和浏览器交互界面。
\end{cvitems}
\end{cventry}

\cvsection{专业技能}
\cvdetail{机器学习}{监督与无监督学习、强化学习、模仿学习、模型评估与泛化。}
\cvdetail{深度学习与机器人}{PyTorch、Transformer、扩散模型、LLaMA 微调、VLA、主动感知、POMDP 与机器人操作。}
\cvdetail{视觉与自动驾驶}{多视图几何、三维重建、目标检测与分割、运动规划、轨迹预测与风险场建模。}
\cvdetail{编程与语言}{C、C++、Python、LaTeX；中文（母语），英语（流利、学术读写）。}
\cvsection{获奖情况}
\cvdetail{2026}{第三届“金灵光杯”中国互联网创新大赛全国二等奖，大学生创新创业成果转化特色赛（团队项目：星火延生）。}
\cvdetail{2026}{腾讯犀牛鸟开源人才培养计划：OpenCloudOS Security Skill issue 实战获奖者，为该题三名获奖者之一。}
\cvdetail{2025}{第五届美团商业分析精英大赛优秀作品奖。}
\cvdetail{2024}{全国大学生数学建模竞赛上海市三等奖。}
\cvdetail{2024}{Water, sanitation, and hygiene for the prevention and care of neglected tropical diseases 入围（日内瓦）。}
\cvdetail{2023}{NVIDIA 全栈 AI 开发工程师技能培训入围。}
\end{document}
